\chapter{转基因、克隆与合成生物学}

\begin{kaichang}
请你想象两样东西。第一样在水果店：一只黄澄澄的夏威夷木瓜，切开后果肉橙红，甜软多汁。第二样在观赏鱼店：一群在紫光灯下发出绿色荧光的小鱼，像一缸游动的霓虹灯管。

这两样东西有什么共同点？它们都是\textbf{被人为改写过 DNA 的生命}。木瓜的基因组里插了一小段病毒基因，小鱼的基因组里插了一个来自水母的发光蛋白基因。

再请你猜两件事：第一，吃下这只木瓜，那段“病毒基因”会不会转到你身上？第二，把荧光小鱼放归野外，会不会毁掉一条河？先写下你的直觉答案。这一章读完，你会发现：这两个问题都不能用“转基因”三个字一刀切地回答——每一个产品，都得单独问、单独查、单独评估。
\end{kaichang}

\section{改写生命这件事，人类干了一万年}

先别急着谈高科技。我们到田野里看几个只用眼睛就能看到的现象。

今天的玉米，穗大粒饱。可它的野生祖先——墨西哥的一种野草“大刍草”——果穗只有小拇指长，上面稀稀拉拉结着十几粒硬得硌牙的种壳。从野草到玉米之间没有发生任何“神奇事件”，只是几千年来，一代代农民做同一件事：\textbf{把长得最合心意的那几株的种子留下来，明年再种}。

再看水果摊。香蕉没有籽，因为它有三套染色体，配对时乱成一锅粥，结不出正常种子；我们吃的其实是靠扦插“克隆”繁殖的植株。大个头的栽培草莓有八套染色体，而森林里的野草莓只有两套。还有狗：吉娃娃和大丹犬长得不像同一个物种，可它们都是狼的后代，全靠人类上万年挑着养。

这些例子说明一件事：\textbf{生物的性状是可以被“改写”的，而且人类早就在改写了}。古人不知道 DNA，他们只知道：挑、留、种，下一代就会朝想要的方向挪一点。他们用的原材料是什么？是每一代个体之间天然存在的差异。

还有一条更“暴力”的路。自然变异的速率毕竟有限，20 世纪的育种家想：能不能人为地多制造一些变异？于是有了诱变育种——把种子送去接受伽马射线照射，或用化学品处理，让 DNA 复制时出更多的错，再从成千上万株“改了不知哪里”的后代里筛出有用的一株。今天全世界用这种方法育成的作物品种超过三千个，你吃的不少大米、大麦、柚子的祖先就进过辐照室。有意思的是：这些品种的基因组被搅得天翻地覆，却很少有人谈“诱变食品”色变。这个现象本身，就值得你在读完本章后回头想一想。

\begin{shiyan}
\textbf{看一看育种的原材料：变异。}材料：两小把绿豆（同一包）、两个浅盘、水、直尺。步骤：把绿豆泡一夜，铺在湿润的纸巾上，两个盘子放在同样的光照和温度下，每天记录发芽时间和芽长，连续观察一周，给每粒豆子编号。观察点：同一包绿豆，哪粒先发芽？一周后芽长相差多少？把芽长按长短排一排，画一张简单的分布图。安全提示：只用清水，泡过生豆子的水不要喝，实验后洗手。你会发现：即便基因背景几乎相同，每粒豆子的表现也有差异——而如果把不同品种的豆子混在一起，差异会更大。育种的全部秘密，就是在这张“差异分布图”上做选择。
\end{shiyan}

\section{把镜头推进：被改写的到底是什么}

现象层只告诉我们“性状会变”。粒子层的问题是：变的到底是什么？

回忆前面的因果链：DNA 序列决定蛋白质，蛋白质在细胞里干活，细胞的整体表现就是你看见的性状。所谓“改写生命”，归根到底只有一件事：\textbf{改变 DNA 的序列}。区别只在于怎么改、改多少、改得准不准。

第 82 章已经认识了这套工具箱：限制性内切酶像“分子剪刀”，连接酶像“胶水”，质粒像“快递车”，CRISPR 像带导航的“文字修改器”。第 83 章又看到这套工具怎样让细菌生产胰岛素、让免疫细胞学会杀癌。这一章我们把同一套工具放到农田、牧场和发酵罐里，并且问一个更尖锐的问题：改了之后，会发生什么？

\section{四种改写方式，一张对比表}

现在引入规则层。市面上争论不休的“育种”“转基因”“基因编辑”，其实是四种不同的 DNA 改写策略：

\begin{table}[htbp]
\centering
\small
\begin{tabular}{@{}p{2.6cm}p{4.6cm}p{2.6cm}p{3.2cm}@{}}
\toprule
方式 & 怎样改变 DNA & 是否定向 & 典型例子 \\
\midrule
传统杂交育种 & 两个亲本整个基因组打乱重组，再人工选择 & 不定向，一次搅动几万个基因 & 杂交水稻、家犬品种 \\
诱变育种 & 用辐射或化学品随机制造大量突变，再筛选 & 不定向，完全随机 & 诱变大麦、太空椒 \\
转基因 & 用分子工具插入一两个特定的外源基因 & 定向引入已知基因 & 抗虫棉、抗病毒木瓜 \\
基因编辑 & 在基因组指定位置改几个碱基 & 最定向，但可能脱靶 & 抗病小麦、高营养番茄 \\
\bottomrule
\end{tabular}
\caption{四种 DNA 改写方式的区别。注意最后一列：它们都只是“手段”，不直接等于“安全”或“危险”。}
\end{table}

请看清楚这张表里最反直觉的一行：传统育种和诱变育种\textbf{搅动的是整个基因组}，而且没人逐个检查被搅动的基因；转基因和基因编辑反而只动一两个位置明确的基因。如果把“改动大小”当作风险指标，排序会和大家想的不一样。

\begin{qiaoliang}
来做一个数量级估算。诱变育种常用的化学诱变剂，能在一株水稻的基因组里留下\textbf{几千个点突变}——这些突变绝大多数没人知道在哪、干了什么，育种家只凭“长得好”把植株留下来，其余几千个突变搭着顺风车一起进了餐桌。而一次典型的转基因操作，引入的是一到两条序列清楚、功能研究过的基因。再给你一个参照：每个人出生时就自带约 60～100 个父母都没有的新发突变（精子、卵子形成时复制的笔误），也就是说，“纯天然、从未突变过的基因组”这种东西，在自然界根本不存在。这不是说改动多少无所谓，而是说：\textbf{风险不在于“改没改过”，而在于“改成什么样”}——这正是下面这条红线。
\end{qiaoliang}

\section{红线：逐个产品评估，不看技术标签}

这一章最重要的一句话，请你记住：\textbf{一个产品的性状和风险，必须逐个评估，不能仅凭“转基因”或“天然”的技术标签来判断。}

逐个评估，具体问三个问题：

\begin{itemize}[itemsep=2pt]
\item \textbf{转入或改动的是什么？}这个基因产生什么蛋白质？这种蛋白质本身有毒性、致敏性吗？抗虫棉转入的 Bt 基因来自苏云金芽孢杆菌，它产生的蛋白质只结合特定昆虫肠道的受体，人的肠道没有这种受体，而且人吃 Bt 蛋白的历史比转基因作物长得多——有机农业拿这种细菌当生物农药喷了几十年。种 Bt 棉之后，棉田化学杀虫剂的喷洒量大幅下降，瓢虫等益虫随之回升；但害虫也不是死靶子——长期单一使用，棉铃虫种群会演化出抗性（这正是第十篇讲的选择压）。于是管理措施跟上：在 Bt 棉田边种一块普通棉花当“庇护所”，让少数没抗性的虫子存活繁殖，稀释抗性基因。你看，连“成功产品”也是一个需要持续运营的动态平衡。
\item \textbf{产生多少、到哪里去？}剂量决定毒性（第 35 章的老规矩）。蛋白质在果实里表达还是在根里表达？会不会随花粉飘到杂草身上？抗除草剂基因如果漂移到近缘野草中，确实可能催生“超级杂草”——这是真实发生过的生态问题，需要田间隔离和管理，而不是一句“绝对安全”能盖过去的。
\item \textbf{换来的性状值不值？}黄金大米转入了两个基因，让胚乳合成 $\beta$-胡萝卜素（维生素 A 的原料），目标是维生素 A 缺乏地区；它经历了二十多年的安全评估和争论才在个别国家获批种植。而历史上第一种转基因食品——延迟成熟的 Flavr Savr 番茄——安全性没有出问题，却因为口味和成本在市场上惨败。\textbf{“安全”与“成功”是两回事，两者都得用事实说。}
\end{itemize}

\begin{zhengju}
科学家怎么知道转基因木瓜既有效、又能通过逐个评估？故事发生在夏威夷。1992 年，木瓜环斑病毒传入夏威夷的主产区，病毒靠蚜虫传播，染病的木瓜果实长满环斑、又小又苦。喷药没用——因为病不在细菌，而在病毒；砍树也没用——蚜虫满天飞。到 1998 年，产量几乎腰斩。

植物病理学家丹尼斯·贡萨尔维斯团队的思路借鉴了“疫苗”的逻辑：把病毒自己外壳蛋白的基因，用第 82 章的工具插进木瓜基因组。植物细胞里出现了少量病毒外壳蛋白，病毒再来入侵时，植物会启动一种类似“沉默同源 RNA”的防御机制，把病毒 RNA 拆掉。但\textbf{想到不等于成立}。他们先做了温室实验：转基因株系接种病毒后确实不发病——可替代解释是“温室环境不同”，于是又把苗种进病毒肆虐的重灾区田间，旁边种普通木瓜作对照。结果：普通植株成片染病，转基因植株健康结果，这是田间对照实验的证据。

接着是安全评估：检测转基因木瓜的蛋白质组成，确认病毒外壳蛋白含量极微（人吃染病但未显症的普通木瓜，吃下的病毒蛋白反而多得多）；比对营养成分，确认除目标性状外与普通木瓜实质等同；评估基因漂流，发现木瓜在夏威夷没有可杂交的近缘野生物种。1998 年，“彩虹木瓜”商业化，夏威夷木瓜产业被救了回来。这个故事的要点不是“转基因万岁”，而是方法：\textbf{有效与否，用对照实验回答；安全与否，逐项检测回答——两个问题都是问这个具体产品，不是问“转基因”这个标签。}
\end{zhengju}

\section{克隆：拷贝的是基因组，不是那个人}

1997 年，一只名叫多莉的绵羊登上全世界报纸的头版。她不是由精子和卵子结合而来，而是科学家把一只成年母羊乳腺细胞的细胞核，移入另一只羊去掉核的卵细胞里，再把这个“拼装细胞”发育成的胚胎植入第三只羊的子宫。多莉的基因组，几乎全部来自那只供核的成年母羊。

镜头推进到粒子层。这件事为什么难？你身上的皮肤细胞和神经细胞带着同一套基因组，却长成完全不同的样子，因为细胞在发育中给大量基因贴上了“关闭”的化学标签（表观遗传修饰）。卵细胞的细胞质是一台神奇的“重启机器”，能把成年细胞核上的标签洗掉，让它重新运行“从受精卵开始”的程序。多莉的成功证明：\textbf{成年体细胞的核里保存着全套遗传信息，没有丢失，只是被调了音}。这也从反面印证了第 74 章的主题：分化不是基因的删减，而是表达方式的改变。

但请注意效率：多莉是 277 次尝试中唯一活到出生的个体。克隆动物流产率高、发育异常多，因为“重启”常常不彻底。多莉本人中年患关节炎、因肺病被安乐死，一度让人担心克隆动物注定早衰。后续研究修正了这个判断：科学家追踪了多莉的基因组“妹妹们”——用同一批细胞克隆出的四只羊，它们健康活到了正常的老年。诚实的说法是：克隆技术会提高发育异常的风险，但\textbf{成功度过发育期的克隆个体可以正常衰老}，早期的“必然早衰”结论被新证据推翻了。

多莉也不是故事的开头。早在 1962 年，约翰·格登就把蝌蚪肠细胞的核移入去核的蛙卵，得到了正常的青蛙——第一次证明成体细胞核里信息齐全。从蛙到羊走了三十五年，因为哺乳动物的卵更小、操作更难。克隆技术今天用在畜牧育种（复制顶尖种畜）和濒危动物保护上；至于克隆宠物，技术上可行，但读完下面这个误解你就会明白，花钱买回的那只“一样的猫”，一样的只是基因组。

\begin{wuqu}
“克隆一个人，就能复制他的记忆和人格。”——不可能。克隆拷贝的只是 DNA 序列，而记忆和人格写在\textbf{神经元的连接方式}里，那是一个人和世界几十年互动的记录，胚胎重编程时被彻底清零。大自然早就做过这个实验：同卵双胞胎就是天然“克隆人”，基因组几乎完全相同，却各自有不同的性格、爱好和人生——因为他们从在子宫里那一刻起，就吃着略有不同的营养、听着不同的声音、做着不同的选择。DNA 是乐谱，不是演奏；克隆给了你同一份乐谱，却给不了同一场演出。
\end{wuqu}

\section{合成生物学：从“改装”走向“设计”}

转基因是往生命这部机器里换一两个零件，合成生物学的野心更大：\textbf{像工程师画电路图一样设计生命}。它有三块招牌。

第一块：最小基因组。一个细胞最少需要多少个基因才能独立活着？克雷格·文特尔团队把支原体近 1000 个基因逐批删减，2016 年造出了只含 473 个基因的人工细胞“辛西娅 3.0”（syn3.0）。最诚实的细节是：这 473 个基因里，约有 149 个\textbf{功能至今未知}——删掉它细胞就死，可没人说得清它干什么。生命这部用了 38 亿年的机器，说明书还有一大半没读懂。这是“我们还不知道”的典型现场。

第二块：人工回路。基因的表达可以被调控，调控就可以连成“电路”：基因 A 的蛋白质抑制基因 B，基因 B 抑制基因 C，基因 C 又回头抑制基因 A——三个基因首尾相咬，就构成一个振荡器，细胞里的蛋白质浓度会像节拍器一样周期性起伏。2000 年，科学家真的在大肠杆菌里搭出了这种“挤压振荡子”。这和收音机里的电子振荡器原理相通，只是零件换成了启动子、基因和蛋白质。下面这幅图是这种思路的人为表示法：

\begin{center}
\begin{tikzpicture}[>=stealth, node distance=1.6cm]
\node[draw, rounded corners, minimum width=2.0cm, minimum height=0.9cm] (A) {基因 A};
\node[draw, rounded corners, minimum width=2.0cm, minimum height=0.9cm, right of=A, node distance=3.2cm] (B) {基因 B};
\node[draw, rounded corners, minimum width=2.0cm, minimum height=0.9cm, right of=B, node distance=3.2cm] (C) {基因 C};
\draw[->, thick] (A) -- (B) node[midway, above]{抑制};
\draw[->, thick] (B) -- (C) node[midway, above]{抑制};
\draw[->, thick] (C.south) to[out=270, in=270] node[midway, below]{抑制} (A.south);
\end{tikzpicture}
\end{center}

三个基因互相抑制，系统就停不下来：A 多时 B 被压低，B 低则 C 升高，C 高又压低 A——如此循环。方框和箭头只是帮助思考的记号，细胞里没有真的电线。

第三块：工程微生物。把通路设计好后，交给微生物这个“活的化工厂”去执行。最漂亮的例子是青蒿素：这种治疟疾的药原本从黄花蒿里提取，产量看天吃饭。科学家把黄花蒿和酵母自身的好几个基因重新组装进酿酒酵母，让这种单细胞真菌在发酵罐里生产青蒿酸，再经一步化学转化得到青蒿素。第 83 章的重组胰岛素是同一个思路：\textbf{细胞是工厂，基因是流水线图纸}。走得更远的团队已经在尝试“从零写染色体”：国际合作的“合成酵母基因组计划”把酵母的染色体一条条用化学方法从头合成出来，再替换掉天然版本，细胞居然照样活着——生命这台机器，正在从“能读、能改”走向“能写”。

\begin{tiaozhan}
为什么工程菌必须“关在”发酵罐里？算一笔账：一个大肠杆菌在理想条件下约 20 分钟分裂一次。1 个细菌经过 24 小时是 72 代，理论上变成 $2^{72} \approx 4.7\times 10^{21}$ 个。每个细菌质量约 $10^{-12}$ 克，总质量约 $4.7\times 10^{9}$ 千克——四百多万吨，比一天内全世界的粮食产量还多。当然，营养会早早耗尽，真实的培养远到不了这个天文数字；但这个估算告诉你：\textbf{微生物的指数增殖能力极其恐怖，一旦带着人工改造的性状逃逸到环境里，没有“撤回键”}。所以合成生物学给工程菌加双保险：让它们依赖自然界不存在的人工营养物（离开工厂就饿死），或装上“杀死开关”（检测到特定信号就自我了断）。跳过这段计算不影响主线，但它解释了后面“生物安全”为什么不是杞人忧天。
\end{tiaozhan}

\section{生物安全：外溢、误用与监管}

最后一类问题超出了一个产品本身，进入了生态和社会层面，仍然要一条条摆事实：

\begin{itemize}[itemsep=2pt]
\item \textbf{生态外溢}。改造过的生物进入环境会怎样？答案照例是“看情况”：抗病毒木瓜在夏威夷没有近缘野草可杂交，外溢风险低；而抗除草剂作物旁边若长着能杂交的近缘杂草，基因漂流就真实发生过。评估的变量是具体的：当地有没有近缘种？性状给不给生存优势？这正是“逐个评估”在生态尺度上的延伸。还有一种新工具把这个问题推到了极限：基因驱动。借助 CRISPR，可以让某个基因以接近 $100\%$ 的概率传给后代（而不是正常的五五开），使“让蚊子不再传播疟疾”这类性状在整个野生种群中迅速扩散。它可能是公共卫生的福音，也可能带来不可预料、且几乎无法撤回的生态后果——所以目前的共识是：先在严格封闭的实验室里研究，野外释放必须逐案经过严格评估和国际协商。
\item \textbf{双重用途}。同一套读写 DNA 的技术，能造疫苗，也可能被用来制造危险的病原体。这不是科幻——合成已知病毒在技术上早已可行。应对办法不是封杀科学（知识本身堵不住），而是监管关键环节：DNA 合成公司筛查订单序列、实验室分级管理、国际公约限制。
\item \textbf{谁受益、谁承担风险}。专利种子和留种权、贫困家庭用不用得起、长期生态监测由谁出钱——这些不是化学问题，但它们决定技术最终长成什么样子。科学给出“能不能”和“会怎样”，社会要回答“该不该”。
\end{itemize}

\section{边界：这套评估框架何时会失灵}

本章给你的核心工具是“逐个产品评估”，但要知道它的边界。第一，基因和性状不是一一对应的说明书关系，一个基因常常影响多个性状（基因多效性），改一处可能带来没料到的副作用——所以短期检测合格不等于可以一劳永逸，商业化后还要持续监测。第二，基因编辑可能脱靶（第 82 章），编辑结果需要逐例测序验证。第三，生态系统太复杂，某些影响要等很多年、放大到很大尺度才显现，任何“长期绝对安全”的承诺都超出了证据。反过来说，“永远不许用”同样是超出证据的断言。诚实的立场是：\textbf{证据走到哪里，政策就走到哪里，边走边测}。

\section{回到那只木瓜和那缸荧光鱼}

现在兑现开场的承诺。

吃下转基因木瓜，那段病毒外壳基因会怎样？和你吃下的米饭、猪肉里的几十亿段基因一样，在消化道里被酶切成核苷酸小分子，绝大部分被分解吸收或排出——你的细胞没有“捡外来 DNA 装到自己基因组里”的机制（第 82 章讲过，连细菌想要吸收外源 DNA 都得人为创造苛刻条件）。何况你吃普通木瓜时，吃下的完整病毒颗粒反而更多。所以第一个问题的答案是：不会，而这个“不会”是从消化与吸收的机制里推出来的，不是从“天然就是好”的标签里喊出来的。

荧光小鱼呢？它转入的水母荧光蛋白对人无害，但它若进入野外河流，可能与本地鱼争夺资源、引入新基因——风险不在荧光，在“外来种”本身，和放生巴西龟是同一类问题（第 86 章讲生态系统时会回到这个网络）。所以第二个问题的答案是：别放生，而这个“别”同样是逐案分析的结果。

你大概已经看出来了：这一章没有给你一份“支持”或“反对”的立场清单，而是给了你一套提问的方法。下一章我们把这套方法用到你自己身上——当一份基因检测报告摆在你面前，告诉你某种疾病“风险升高 $30\%$”时，你该怎样读懂它、又不被它吓住？

\section*{练一练}

\begin{sikaoti}
\item 画一张信息流图：从“抗病毒木瓜里的病毒外壳基因”出发，画出它经过转录、翻译到“植物抵抗病毒”的完整链条，并在每个箭头上注明发生了什么事。图旁注明：这是信息流的表示法，不是细胞的写实图。
\item 参考本章的对比表，用自己的话向一位坚信“杂交育种天然安全、转基因必然危险”的长辈解释：为什么“改动的基因数量”不是风险的可靠指标？请用上诱变育种的例子和突变数量级。
\item 假设有人提议“把抗除草剂基因转到油菜里，推广到你所在地区种植”。请列出你会要求评估的至少四个具体问题（提示：从基因、蛋白质、剂量、花粉漂流、本地近缘植物几个角度想）。
\item 多莉羊的基因组全部来自供核母羊，但她出生后吃的草、得的病、接受的人为照料都是自己的经历。据此画一幅“同一份乐谱、两场不同演出”的示意图，说明基因组相同而性状可以不同，并指出这种差异来自哪些环节。
\item 一个大肠杆菌 20 分钟分裂一次，从 1 个开始，10 小时后理论上有多少个？（用 $2$ 的幂表示，再换算成以 10 为底的数量级。）如果营养只允许总质量达到 1 克，增殖会在什么时候“刹车”？这说明发酵罐和自然环境分别靠什么限制微生物？
\item 本章说“同卵双胞胎是天然克隆人”。请设计一个观察思路，用来检验“基因组相同是否意味着性格相同”——你会比较谁和谁？有哪些混杂因素需要先排除？
\end{sikaoti}
